\documentclass[a4paper,11pt]{article}

\usepackage[utf8]{inputenc}
\usepackage[T1]{fontenc}
\usepackage[french]{babel}

\usepackage{amsmath}
\usepackage{amssymb}
\usepackage{graphicx}
\usepackage{geometry}
\usepackage{hyperref}
\usepackage{listings}
\usepackage{xcolor}
\usepackage{booktabs}
\usepackage{tikz}

\geometry{margin=2.5cm}

\lstset{
language=Python,
basicstyle=\ttfamily\small,
keywordstyle=\color{blue},
commentstyle=\color{gray},
stringstyle=\color{red},
frame=single,
breaklines=true,
showstringspaces=false
}

\title{\textbf{Calibration de la position du Tool Center Point (TCP)}}

\date{\today}

\begin{document}

\maketitle

\begin{abstract}

La précision de positionnement d’un outil dépend fortement de la connaissance exacte de sa géométrie par rapport au système mécanique qui le porte.  
Le \textit{Tool Center Point} (TCP) correspond au point de référence utilisé pour décrire la position réelle de l’outil dans l’espace.

Ce rapport présente une méthode de calibration du TCP basée sur la collecte de plusieurs poses autour d’un point fixe.
Les données sont ensuite exploitées à l’aide d’une résolution par moindres carrés afin d’estimer l’offset du TCP dans le repère de la bride.  
Finalement, différentes métriques permettent de vérifier la validité des mesures ainsi que la qualité de la calibration obtenue.

La méthode de pivot utilisée permet uniquement d’estimer la position du TCP.
L’orientation de l’outil ne peut pas être déterminée avec cette approche car aucune contrainte géométrique ne permet de l’observer.

\end{abstract}

\tableofcontents
\newpage

\section{Introduction}

Dans de nombreuses applications de manipulation ou d’usinage, les trajectoires sont définies par rapport à un point spécifique de l’outil.  
Ce point, appelé \textit{Tool Center Point} (TCP), correspond généralement à l’extrémité fonctionnelle de l’outil.

La position du TCP n’est pas directement connue par le système mécanique.  
Elle doit être définie comme un décalage (\textit{offset}) par rapport à la bride de montage de l’outil.

Une estimation incorrecte de cet offset peut entraîner :

\begin{itemize}
\item des erreurs de positionnement
\item des trajectoires imprécises
\item une mauvaise répétabilité des opérations
\end{itemize}

La calibration du TCP consiste donc à déterminer avec précision la position de ce point dans le repère de la bride.

\section{Définition des repères}

La calibration repose sur trois repères principaux :

\begin{itemize}
\item repère de base
\item repère de la bride (flange)
\item repère outil TCP
\end{itemize}

\begin{figure}[h]
\centering
\begin{tikzpicture}[scale=1.2]

\draw[->] (0,0) -- (2,0) node[right]{$x_{base}$};
\draw[->] (0,0) -- (0,2) node[above]{$y_{base}$};

\draw[->,blue] (3,1) -- (4,1) node[right]{$x_{flange}$};
\draw[->,blue] (3,1) -- (3,2) node[above]{$y_{flange}$};

\filldraw (3,1) circle (2pt) node[below] {Bride};

\filldraw (4,2) circle (2pt) node[right] {TCP};

\draw[dashed] (3,1) -- (4,2);

\end{tikzpicture}

\caption{Relation entre la bride et le TCP}
\end{figure}

\section{Principe de la calibration du TCP}

La méthode utilisée repose sur une propriété géométrique simple :

\begin{center}
"le TCP correspond à un point fixe dans l’espace lorsque l’outil pivote autour de celui-ci."
\end{center}

Lors de la calibration, l’extrémité de l’outil est maintenue sur un point fixe tandis que l’orientation (angle d'approche) du système change.

\begin{figure}[h]
\centering
\begin{tikzpicture}[scale=1]

\filldraw (0,0) circle (3pt) node[below] {Point pivot};

\draw[thick] (0,0) -- (1,1);
\draw[thick] (0,0) -- (-1,1);
\draw[thick] (0,0) -- (2,0);

\end{tikzpicture}

\caption{Principe du pivot utilisé pour la calibration}
\end{figure}

Chaque mesure fournit :

\begin{itemize}
\item la position de la bride
\item son orientation
\end{itemize}

Ces informations permettent de reconstruire la position du TCP par rapport à la bride.

\section{Collecte des données}

\subsection{Principe}

La calibration nécessite plusieurs poses différentes autour du même point pivot.

Chaque pose est représentée par :

\[
[x,y,z,r_x,r_y,r_z]
\]

où :

\begin{itemize}
\item $(x,y,z)$ représente la position de la bride
\item $(r_x,r_y,r_z)$ représente le vecteur rotation (axis-angle)
\end{itemize}

Ces valeurs permettent de reconstruire la transformation homogène :

\[
T =
\begin{bmatrix}
R & p \\
0 & 1
\end{bmatrix}
\]

avec :

\begin{itemize}
\item $R$ la matrice de rotation
\item $p$ le vecteur de translation
\end{itemize}

\subsection{Procédure}

La collecte des données se déroule selon les étapes suivantes :

\begin{enumerate}
\item placer la pointe de l’outil sur un point fixe
\item modifier l’orientation du système
\item enregistrer la pose
\item répéter l’opération plusieurs fois
\end{enumerate}

Un minimum de six poses est nécessaire afin d’obtenir une solution stable, mais un nombre plus élevé de mesures améliore la robustesse du calcul.

\section{Validation des données collectées}

Avant d'effectuer la calibration, il est important de vérifier la qualité des mesures.

\subsection{Cohérence du point pivot}

À partir d’un TCP de référence $
r=
\begin{bmatrix}
    x,y,z
\end{bmatrix}
$ comme un ancienne calibration, on peut estimer la position du point pivot dans le repère global :

\[
c_i = p_i + R_i r
\]

avec :

\begin{itemize}
\item $R_i$ la matrice de rotation
\item $p_i$ le vecteur de translation
\item $c_i$ la position TCP théorique 
\end{itemize}

Si les mesures sont cohérentes, les points $c_i$ doivent être très proches les uns des autres.

La dispersion est définie par :

\[
spread = \max ||c_i - \bar{c}||
\]

Une faible valeur indique que les mesures sont compatibles avec l’hypothèse d’un point fixe.

\subsection{Amplitude du mouvement}

Pour éviter un problème mal conditionné, la bride doit avoir suffisamment bougé :

\[
motion = \max ||p_i - p_0||
\]

\subsection{Variation d’orientation}

Une variation importante d’orientation améliore la précision de la calibration.

L’angle de rotation peut être estimé à partir de la trace de la matrice de rotation :

\[
\theta_i = \cos^{-1}\left(\frac{trace(R_i)-1}{2}\right)
\]

\[
\Delta\theta = \max(\theta_i) - \min(\theta_i)
\]

\section{Calcul du TCP}

Le TCP est défini comme un vecteur $t$ dans le repère de la bride.

Pour chaque mesure :

\[
c = p_i + R_i t
\]

où $c$ correspond au point pivot dans le repère global.

Cette relation peut être réécrite :

\[
R_i t - c = -p_i
\]

En empilant toutes les équations on obtient :

\[
A_i x_i = b_i
\]

avec

\[
x =
\begin{bmatrix}
t \\
c
\end{bmatrix}
\]

\[
A =
\begin{bmatrix}
R & I
\end{bmatrix}
\]

Le système est résolu par une méthode de \textbf{moindres carrés} :

\[
x = (A^TA)^{-1}A^Tb
\]

Cette solution fournit :

\begin{itemize}
\item $t$ la position $
\begin{bmatrix}
    x,y,z
\end{bmatrix}
$ du TCP dans le repère de la bride
\item $c$ la position du point pivot dans le repère global
\end{itemize}

\newpage
\section{Validation de la calibration}

Une fois le nouvel offset déterminé, il est nécessaire de vérifier que la calibration est correcte.

\subsection{Test de rotation}

Le principe consiste à effectuer des rotations autour du TCP estimé :

\begin{itemize}
\item rotation autour de l’axe $x$
\item rotation autour de l’axe $y$
\item rotation autour de l’axe $z$
\end{itemize}

dans les deux directions.

\subsection{Observation}

Si la calibration est correcte :

\begin{itemize}
\item le TCP reste fixe
\item seule l’orientation change
\end{itemize}

Si un déplacement du point est observé, cela signifie que l’offset TCP est incorrect.

\section{Implémentation}

Algorithme:

\begin{enumerate}

\item Collecter plusieurs poses
\item Convertir les poses en matrices homogènes
\item Construire le système linéaire
\item Résoudre par moindres carrés
\item Extraire l’offset TCP
\item Valider par rotation autour du TCP

\end{enumerate}

\section{Conclusion}

La calibration du TCP est une étape essentielle pour garantir la précision des mouvements basés sur un point outil.

La méthode présentée repose sur :

\begin{itemize}
\item la collecte de plusieurs poses autour d’un point fixe
\item la résolution d’un système linéaire par moindres carrés
\item la validation expérimentale de la calibration
\end{itemize}

Cette approche permet d’obtenir une estimation robuste du TCP tout en restant simple à mettre en œuvre.

\newpage
\appendix

\section{Code Python}

\subsection{Exemple simplifié de calcul TCP}
\begin{lstlisting}
import numpy as np

def calibrate_tcp(Ts):
    '''
    Ts: List of homogeneous transformation matrices representing
        the pose of the robot flange for each measurement.

    return: positional TCP offset (x,y,z)
    '''
    N = len(Ts)

    A = np.zeros((3*N,6))
    b = np.zeros((3*N))

    for i,T in enumerate(Ts):

        R = T[:3,:3]
        p = T[:3,3]

        A[3*i:3*i+3,0:3] = R
        A[3*i:3*i+3,3:6] = -np.eye(3)
        b[3*i:3*i+3] = -p

    x,_,_,_ = np.linalg.lstsq(A,b,rcond=None)

    return x[:3]

\end{lstlisting}

\end{document}